% =====================================================================
% EasyNet Ability Owner — Truth Table & Default Catalogue
% =====================================================================
%
% Companion to:
%   docs/rfc/AXON-RFC-001-plan-v4.1.5.md  (URA + AXIOM seven-tuple)
%
% Purpose: define, exhaustively, how every default-registered ability
% maps to the AXIOM Invocation envelope, so future contributors can
% verify owner stamping, render-layer projection, and federation
% advertise behaviour without re-deriving the analysis from scratch.
%
% Author: 凉冰 (architect)
% Date:   2026-05-05
%
% Compile: xelatex (twice for hyperref bookmarks).

\documentclass[10pt,a4paper]{article}

\usepackage[utf8]{inputenc}
\usepackage[T1]{fontenc}
% xelatex picks up the system PingFang / Songti fonts for CJK fallback.
% Wrapping the few CJK glyphs (the architect's name 凉冰) in xeCJK
% lets them render without forcing the whole document onto a Chinese
% font.
\usepackage{fontspec}
\usepackage{xeCJK}
\setCJKmainfont[BoldFont={PingFang SC Semibold}]{PingFang SC}

\usepackage[a4paper,margin=2.2cm,top=2.4cm,bottom=2.4cm]{geometry}
\usepackage{lmodern}
\usepackage{microtype}
\usepackage{xcolor}
\usepackage{array}
\usepackage{longtable}
\usepackage{booktabs}
\usepackage{tabularx}
\usepackage{ltablex}
\keepXColumns
\usepackage{enumitem}
\usepackage{fancyvrb}
\usepackage{titlesec}
\usepackage{caption}
\usepackage{hyperref}

% ---------- Hyperref styling ------------------------------------------
\hypersetup{%
  colorlinks   = true,
  linkcolor    = black,
  citecolor    = black,
  urlcolor     = black!70,
  pdftitle     = {EasyNet Ability Owner Truth Table},
  pdfauthor    = {Liangbing (architect)},
  pdfsubject   = {AXIOM seven-tuple owner mapping for the default ability catalogue},
}

% ---------- Section headings -------------------------------------------
\titleformat{\section}{\large\bfseries}{\thesection}{0.6em}{}
\titleformat{\subsection}{\bfseries}{\thesubsection}{0.6em}{}
\setlength{\parskip}{0.45em}
\setlength{\parindent}{0pt}

% ---------- Code style -------------------------------------------------
\renewcommand{\ttdefault}{lmtt}
\newcommand{\code}[1]{\texttt{\small #1}}

% ---------- Break-friendly typewriter -----------------------------------
% Long URA / fully-qualified-name strings cannot break inside a normal
% \texttt cell. The `url` package parses verbatim and inserts soft
% breaks at the listed separators. \codew is a tt-styled \url variant;
% the visual matches \code when no break fires.
\usepackage{url}
\makeatletter
\g@addto@macro{\UrlBreaks}{\do\.\do\:\do\/\do\_\do\-\do\<\do\>}
\makeatother
\DeclareUrlCommand\codew{\urlstyle{tt}\renewcommand{\UrlFont}{\ttfamily\small}}

% ---------- Custom column types ---------------------------------------
\newcolumntype{L}[1]{>{\raggedright\arraybackslash}p{#1}}
\newcolumntype{C}[1]{>{\centering\arraybackslash}p{#1}}

% ---------- Table macros ----------------------------------------------
\definecolor{rowdim}{gray}{0.94}
\renewcommand{\arraystretch}{1.18}

% =====================================================================
\begin{document}

\begin{center}
  {\LARGE\bfseries EasyNet Ability Owner --- Truth Table}\\[0.4em]
  {\large\textit{AXIOM Invocation envelope projection for the default catalogue}}\\[0.3em]
  \rule{0.55\textwidth}{0.4pt}\\[0.4em]
  {\small Author: \emph{凉冰} (architect) \quad Date: 2026-05-05 \quad Status: \textsc{Specification}}\\[0.2em]
\end{center}

\vspace{0.5em}

\section*{Why this document exists}

This is a flat-fact reference: a single page-pair that answers, for every
default-registered ability, three questions:

\begin{enumerate}[leftmargin=1.6em,itemsep=0.1em,topsep=0.1em]
  \item \emph{What URA names its owner} (the \code{owner\_ura}
        field on its descriptor).
  \item \emph{Where it lands in the AXIOM Invocation envelope}
        (\code{caller} / \code{callee} / \code{subject} slots).
  \item \emph{Why EasyNet ships it by default} (the protocol assumption
        that breaks if you delete it).
\end{enumerate}

The CLI render layer, \code{meta.list\_abilities} synth path, and
federation advertise pipeline all derive owner mapping independently.
The truth table is the contract those derivations must agree with.

% =====================================================================
\section{URA grammar (AXON-RFC-001 v4.1.5 §A.URA)}

Every value referenced in this document conforms to the v4.1.5 URA
grammar. The grammar is closed: the parser
(\codew{src/core/ura/mod.rs::parse\_ura}) rejects any string outside the six
shapes below. Other "URA-shaped" tokens that have appeared in practice
(\codew{legacy self alias}, \codew{<unjoined>}, \code{agent://self}) are
\emph{not} URAs and must not propagate past the registration boundary.

\subsection*{Ontology}

\begin{tabularx}{\linewidth}{l X}
  \toprule
  \textbf{Relation} & \textbf{Meaning} \\
  \midrule
  realm \emph{owns} hub        & A realm is identified by its hub URA; one hub per realm. \\
  user \emph{owns} agent       & Ownership stays with the user across device migrations. \\
  device \emph{hosts} agent    & Device is the runtime placement, not the identity. \\
  agent \emph{exposes} ability & Verbs are mounted under their owning agent. \\
  resource \emph{is} user-owned concrete substrate & Files, processes, ptys, sockets — typed by namespace. \\
  \bottomrule
\end{tabularx}

\subsection*{The six URA shapes}

\begin{tabularx}{\linewidth}{l l X}
  \toprule
  \textbf{Kind} & \textbf{Wire shape} & \textbf{Tail invariants} \\
  \midrule
  \code{user}     & \codew{easynet:///r/<realm>/user/<user-uuid>}                           & user-uuid is bare (no \code{.}, no \code{/}) \\
  \code{device}   & \codew{easynet:///r/<realm>/device/<device-uuid>}                       & device-uuid is bare \\
  \code{agent}    & \codew{easynet:///r/<realm>/agent/<user-uuid>.<agent-id>}               & exactly one \code{.} in tail; both segments non-empty; agent-id has no \code{.} \\
  \code{ability}  & \codew{easynet:///r/<realm>/ability/<user-uuid>.<agent-id>.<ability-id>} & exactly two leading \code{.}s; ability-id may itself contain \code{.}s (e.g.\ \code{fs.read}) \\
  \code{hub}      & \codew{easynet:///r/<realm>/hub}                                        & realm-singleton; no tail; no sub-id; extra tail segments are invalid \\
  \code{resource} & \codew{easynet:///r/<realm>/resource/<user-uuid>/<ns>/<path>}           & \code{<ns>}~$\in$~\{\code{fs}, \code{process}, \code{pty}, \code{shell}, \code{http}\}; v4.1.5 also accepts \codew{resource/<user>.<owner-tail>/<path>} for content-addressed owners \\
  \bottomrule
\end{tabularx}

\subsection*{Three properties used throughout this document}

\begin{description}[leftmargin=1.4em,style=nextline]
  \item[\textbf{P1} --- URAs are self-describing] Any URA, in
  isolation, names exactly one actor in the seven-tuple ontology.
  Nothing about who \emph{calls} the URA changes what it names. A
  string that requires caller context to resolve (\codew{legacy self alias}, the
  literal placeholder \code{self}, an \code{agent://self} URN) is
  therefore not a URA, and must not appear in any persisted descriptor,
  advertise payload, or receipt field.

  \item[\textbf{P2} --- Owner is named at registration, not derived
  from the name] The descriptor's \code{owner\_ura} is the
  wire form of "who advertises this ability". It is fixed at
  \code{register\_rpc} time, by the registration site, against an
  explicit owner URA. It is \emph{not} inferred from
  \code{name.starts\_with(\dots)} prefix sniffing. Every consumer
  reads the stored value verbatim.

  \item[\textbf{P3} --- \code{callee} URA carries the actor, not the
  verb] The \code{ability} URA kind exists only at the publish layer
  (descriptor, advertise). It never appears in an Invocation
  envelope's \code{callee} slot. Per AXON-RFC-001 v4.1.5 §9
  (AXIOM seven-tuple), \code{callee}~$\in$~\code{\{hub, device, agent\}}. The
  verb itself is in \code{function\_name}, which is a string, not a
  URA.
\end{description}

% =====================================================================
\section{The AXIOM Invocation envelope}

Per AXON-RFC-001 v4.1.5 §9 the Invocation envelope carries seven
URA-typed bindings; only the four that govern owner mapping are
reproduced here.

\begin{tabularx}{\linewidth}{l l X}
  \toprule
  \textbf{Field} & \textbf{Allowed URA kind} & \textbf{Semantics} \\
  \midrule
  \code{caller}   & \code{hub} / \code{device} / \code{agent} & Who initiates the call. Required by Invariant~1. \\
  \code{callee}   & \code{hub} / \code{device} / \code{agent} & Which actor runs the handler. \emph{Never} \code{ability} or \code{user}. \\
  \code{subject}  & any URA kind, may equal \code{callee} & Resource the call acts on. \code{None} is allowed only for meta-style abilities (\textsc{INV-META-SUBJECT-EXEMPT}); resource handlers must reject \code{None}. \\
  \code{function\_name} & string, not a URA & The ability name (\codew{fs.read}, \codew{<owner>.chat}, \codew{hub.openai.chat_completions}, \dots). \\
  \bottomrule
\end{tabularx}

% =====================================================================
\section{Truth table: name shape \texorpdfstring{$\to$}{->} owner kind}

Every default-registered ability falls into one of the following
name shapes. For each shape the table fixes the canonical
\code{owner\_ura}, the resulting \code{envelope.callee}, and
the default \code{envelope.subject} convention. The \emph{Where}
column points to the registration call site, which is the
authoritative source for the owner per~P2.

\renewcommand{\arraystretch}{1.32}
\setlength{\tabcolsep}{4pt}
\sloppy
\emergencystretch=3em
\begin{longtable}{@{}p{0.35cm} L{4.0cm} L{3.9cm} L{3.6cm} L{3.4cm}@{}}
  \toprule
  \textbf{\#} & \textbf{Name shape} &
  \textbf{Canonical owner URA} &
  \textbf{Default \code{subject}} &
  \textbf{Where} \\
  \midrule
  \endhead

  1 & \codew{device.fs.*}, \codew{device.http.*}, \codew{device.shell.*}, \codew{device.process.*} &
      \code{device/<id>} &
      resource URA, or \code{=callee} for meta verbs &
      AXIOM Tier~2.5 baseline-locomotion-v1 profile \\

  2 & \codew{device.camera.*}, \codew{device.mic.*}, \codew{device.screen.*}, \codew{device.speaker.*} &
      \code{device/<id>} &
      \emph{required} resource URA; handler rejects \code{None} &
      \codew{profiles::device} (media) \\

  3 & \codew{device.fleet.*}, \codew{device.observe.*}, \codew{device.admin.*}, \codew{device.discuss.*}, \codew{device.loop.*}, \codew{device.schedule.*}, \codew{device.mission.*}, \codew{device.meta.*}, \codew{device.ability.*}, \codew{device.policy.*}, \codew{device.consent.*}, \codew{device.skill.*}, \codew{device.mcp.*}, \codew{device.a2a.*} &
      \code{device/<id>} &
      \code{=callee} (meta) &
      \codew{profiles::device}; runtime \codew{register_rpc} \\

  4 & \codew{<agent>.discover}, \codew{<agent>.invoke}, \codew{<agent>.api_key.*}\newline
      (one row per mounted sub-agent; \codew{<agent>}~$\in$~\{\code{claude}, \code{codex}, \code{web-builder}, \dots\}) &
      \code{agent/<u>.<agent>} &
      \code{=callee} (meta) &
      \codew{discover_ability::register_for_agent}\newline
      \codew{invoke_ability::register_for_agent}\newline
      \codew{api_key_ability::register} \\

  5 & \codew{device.keyring.*}, \codew{session.open}, \codew{identity.register_pubkey} &
      \code{device/<id>} &
      \code{=callee} (RFC-002 §3.3) &
      \codew{keyring::abilities::register_for_owner}\newline
      \codew{axon_serve::admission_facade} \\

  6 & \codew{<agent>.chat} &
      \code{agent/<u>.<agent>} &
      \code{=callee} (chat is meta) &
      \codew{chat_ability::register} \\

  7 & \codew{<agent>.<custom-verb>}\newline
      (e.g.\ \codew{web-builder.todo_add_task}) &
      \code{agent/<u>.<agent>} &
      task or resource URA &
      per-agent profile \code{register} \\

  8 & \code{consent.*}, \code{policy.*} (handler form), \code{mcp.bridge.*}, \code{mcp.client.*}, \code{a2a.bridge.*}, \code{a2a.client.*}, \code{skill.*} &
      \code{agent/<u>.<helper-id>} for hosted helper agents (e.g.\ \codew{consent-default-0}); \code{device/<id>} for daemon-side adapters &
      varies; resource URAs for outbound adapters &
      \codew{profiles::consent}; \codew{profiles::policy}; \codew{profiles::mcp} \\

  9 & \code{voice.*}, \codew{hub.openai.*}, \codew{hub.<namespace>.*} &
      \code{hub} (\codew{easynet:///r/<r>/hub}, realm-singleton) &
      \code{=callee} or per-verb resource URA &
      \codew{voice::register}; \codew{media::register}; \codew{openai_compat_ability::register} \\

  10 & \code{<user>.pages.*}, \codew{<user>.<project>.page.fetch}, \code{<user>.files.*} &
      \code{agent/<u>.<owner-tail>} or \code{user/<u>} depending on profile &
      resource URA (\codew{resource/<u>.<owner-tail>/...}) &
      \codew{pages::register}; \codew{files::register} \\

  11 & \codew{<user>.api_key.*}\newline
       (\code{<user>} is the canonical user-id from credentials, not a placeholder) &
      \code{user/<user>} &
      \code{=callee} (meta) &
      \codew{api_key_ability::register} \\

  \bottomrule
\end{longtable}

\subsection*{Notes on the table}

\begin{description}[leftmargin=1.4em,style=nextline]
  \item[Row 4 is per-sub-agent] The catalogue contains one
  \codew{discover}/\codew{invoke}/\codew{api_key.*} row \emph{per
  mounted sub-agent}. With sub-agents \code{claude} and \code{codex}
  on the same device, the catalogue carries
  \code{claude.discover}+\code{codex.discover}+\dots, each with its
  own owner URA. There is no shared "self" entry.

  \item[Row 5 is exactly one row] Daemon-bundle abilities
  (\code{device.keyring.*}, \code{device.session},
  \code{device.register\_device\_pubkey}) are mounted once per daemon,
  with \code{owner\_ura = device URA}. They are not per-agent.

  \item[Row 8 has two sub-cases] An adapter such as \code{consent.decide}
  can be hosted either as a sub-agent (\codew{consent-default-0}, owner
  = \code{agent/<u>.consent-default-0}) or as a daemon-side bridge
  (\code{mcp.bridge.list\_tools}, owner = \code{device/<id>}). The
  catalogue records the sub-case at registration time --- see~P2.

  \item[Row 9 owner is the realm hub] \codew{hub.*} handlers physically
  run on the device daemon (the hub forwards via federation), but the
  ability is conceptually a hub-tier projection. The descriptor must
  say so. Synth derives the owner URA via
  \codew{hub_ura(credentials.tenant_id)} --- not from the device URA.

  \item[Row 11 \code{<user>} is canonical] The \code{<user>} slot is
  the canonical user-id from \code{credentials.username}; placeholder
  defaults are not part of the spec.
\end{description}

% =====================================================================
\section{CLI render projection}

The \code{easynet ability list} table projects the
\code{owner\_ura} into the columns
\textsc{device}~/~\textsc{agent}~/~\textsc{user}. Each column reflects
a slot of the owner URA's kind; columns the owner does not name stay
as a literal dash. Mixing kinds in one row (the
\code{device=99e59cc... agent=hub} bug from this PR's audit) is a
renderer-level violation of P3 and must be rejected by tests.

\begin{tabularx}{\linewidth}{l l l l X}
  \toprule
  \textbf{Owner URA kind} & \textsc{device} & \textsc{agent} & \textsc{user} & \textbf{Notes} \\
  \midrule
  \code{device/<id>}            & \code{<id>} & \code{-}    & \code{-}    & rows~1, 2, 3, 5 \\
  \code{agent/<u>.<a>}          & \code{-}    & \code{<a>}  & \code{<u>}  & rows~4, 6, 7, 8a, 10a \\
  \code{hub}                    & \code{-}    & \code{hub}  & \code{-}    & row~9 \\
  \code{user/<u>}               & \code{-}    & \code{-}    & \code{<u>}  & row~11 (canonical) \\
  parse failure / unknown       & \code{-}    & \code{!err} & \code{-}    & catalogue corruption signal; do \emph{not} silently dash \\
  \bottomrule
\end{tabularx}

% =====================================================================
\section{Default catalogue --- abilities and why each exists}

The 121 abilities below are what \code{easynet ability list} returns
on a freshly-paired daemon (realm \codew{easynet.run}, user \code{test},
device \code{99e59ccd-...debf73d}, with one hosted sub-agent under the
\code{web-builder} profile). For each ability we state \emph{why
EasyNet must ship it by default} --- the protocol assumption that
breaks if you delete the row.

% ---------- Hub (2) ----------
\subsection{Hub-owned (owner = \code{hub})}

\paragraph{\texttt{\small hub.openai.chat\_completions}, \texttt{\small hub.openai.list\_models}.}
RFC-006-C v0.1 OpenAI compatibility surface. Without these, an OpenAI
SDK pointed at an EasyNet daemon cannot do non-streaming chat or model
enumeration --- the entire "drop-in OpenAI replacement" story fails at
the wire. The verbs are hub-rooted (not device-rooted) because the
realm hub is the federation-visible endpoint that external SDKs dial;
the device daemon executes as the hub's local-execution proxy. The
legacy on-wire form is \codew{01HUB.openai.*}; the rename to
\codew{hub.*} ships with RFC-001 v4.1.6.

% ---------- Device baseline (Tier 2.5) ----------
\subsection{Device-owned --- AXIOM Tier 2.5 baseline-locomotion-v1 (under \texttt{\small device.*})}

These four namespaces are the AXIOM \emph{baseline locomotion} profile.
Every Tier-2.5-conformant device ships them; the 8-stage destructive-
intent / argv-parsing / receipt-redaction pipeline they share is the
\emph{reason} EasyNet can claim "an LLM can move on this substrate
without bypassing audit". Removing any of them collapses the audit
contract.

\begin{description}[leftmargin=1.4em,style=nextline,itemsep=0.15em]
  \item[\codew{fs.edit}, \codew{fs.list}, \codew{fs.read}, \codew{fs.write}]
  Filesystem locomotion. \codew{fs.edit} is the surgical
  string-replace primitive that lets agents diff-apply rather than
  full-rewrite (audit-trail compactness).
  \item[\codew{http.request}] One outbound HTTP request per
  invocation. The receipt captures status, headers, body bytes for
  the audit pipeline; without this primitive every web-fetching
  agent has to escape into shell and the audit trail is forfeit.
  \item[\codew{shell.run}] Bash command string with the full
  AXIOM Tier-2.5 8-stage security pipeline. The destructive-intent
  classifier and shell tokeniser live here.
  \item[\codew{process.exec}] OS-level argv (no shell interpretation)
  --- the surface for tools that need predictable argument passing
  without the shell metacharacter risk surface. Pairs with
  \codew{shell.run} as the safe-default counterpart.
\end{description}

\subsection{Device-owned --- media resources (under \texttt{\small device.*}, subject required)}

The four media namespaces project hardware resources into the
ability surface. They are the \emph{only} way an agent on this device
can see/hear/speak; no fallback path exists. Subject (a resource URA
naming the camera/mic/screen/speaker handle) is mandatory --- handlers
return \code{resource\_not\_found} on \code{None} per
\textsc{INV-SUBJECT-ENVELOPE}.

\begin{description}[leftmargin=1.4em,style=nextline,itemsep=0.15em]
  \item[\codew{camera.snapshot}, \codew{camera.subscribe}] Single
  still vs.\ pushed video stream from a camera resource.
  \item[\codew{mic.subscribe}] Pushed audio stream from a microphone
  resource.
  \item[\codew{screen.snapshot}, \codew{screen.subscribe}] Single
  still vs.\ pushed video stream of a display target.
  \item[\codew{speaker.publish}] Outbound audio frames to a speaker
  resource. A caller may route Hub-owned \codew{voice.subscribe} TTS
  output into this Device-owned sink through a separate Invocation.
\end{description}

\subsection{Device-owned --- fleet operational surface (under \texttt{\small device.fleet.*})}

The fleet namespace exposes \emph{cross-device operator verbs through
the same Invoke surface as everything else}. Without them, operators
have to ssh into each device --- which forfeits AXIOM auditability.
Per RFC §18 the device profile owns these.

\begin{description}[leftmargin=1.4em,style=nextline,itemsep=0.15em]
  \item[Node lifecycle:]
  \codew{device.fleet.list_nodes}, \codew{device.fleet.describe_node},
  \codew{device.fleet.remove_node}.
  Operator-facing fleet topology view.

  \item[Ability deployment:]
  \codew{device.fleet.deploy_ability}, \codew{device.fleet.uninstall_ability},
  \codew{device.fleet.list_abilities}.
  Cross-device ability publish + uninstall via the Invoke surface
  (mirrors \code{easynet ability deploy}).

  \item[Agent lifecycle:]
  \codew{device.fleet.start_agent}, \codew{device.fleet.stop_agent},
  \codew{device.fleet.list_agents}.
  Sub-agent registry mutations through Invoke (mirrors
  \code{easynet agent add} / \code{remove}).

  \item[Skill management:]
  \codew{device.fleet.skill_install}, \codew{device.fleet.skill_remove},
  \codew{device.fleet.skill_upgrade}.
  Marketplace skill ops on a target device's
  \code{<agent-root>/skills/} dir.

  \item[Sessions (interactive PTY):]
  \codew{device.fleet.session_create}, \codew{device.fleet.session_attach},
  \codew{device.fleet.session_input}, \codew{device.fleet.session_read},
  \codew{device.fleet.session_resize}, \codew{device.fleet.session_close},
  \codew{device.fleet.list_sessions}, \codew{device.fleet.attach_session}.
  PTY-aware interactive shell over Invoke; pairs with
  \codew{device.fleet.exec_remote} (one-shot) for non-interactive cases.
  The \codew{pty_session_*} variants are explicit aliases retained
  for legacy callers.

  \item[Direct ops:]
  \codew{device.fleet.exec_remote}, \codew{device.fleet.file_transfer}.
  One-shot remote command + bidirectional file copy. Without these
  the operator has no way to land a file or check a status without
  out-of-band access.
\end{description}

\subsection{Device-owned --- observability, admin, discovery aliases (under \texttt{\small device.*})}

\begin{description}[leftmargin=1.4em,style=nextline,itemsep=0.15em]
  \item[\codew{admin.status}] Operator-facing component snapshot
  (build version, membership state, ability-registry size). Pairs
  with the two \code{observe.*} verbs as the daemon's three-tier
  health surface.
  \item[\codew{device.describe}] Device-profile self-description.
  Cross-device addressing helper (lets a peer ask "are you who I
  think you are" without pulling the full ability catalogue).
  \item[\codew{observe.health}] Round-trip smoke ability. Returns
  the caller's args plus a daemon timestamp. The minimal proof of
  liveness; every CI sanity check uses this.
  \item[\codew{observe.network_health}] Live network posture: local
  reachability, realm-directory reachability, peer-hub status. The
  diagnostic that distinguishes "daemon down" from "daemon up but
  hub unreachable".
  \item[\codew{mission.track}, \codew{mission.cancel}]
  Mission lifecycle introspection + termination. \code{track}
  reads the persisted run dir of a prior \codew{mission.run}
  invocation; \code{cancel} flips an in-flight mission to
  \texttt{cancelled}. RFC-001 v4.1.7 M2 retired the parallel
  \codew{device.easynet.*} alias surface that previously mirrored
  \code{run}/\code{track}/\code{cancel} (and \code{discover} as an
  alias for \codew{device.meta.list\_abilities}); per Q2 of the
  migration plan aliases are protocol entropy and the LLM corpus
  reads the canonical names.
\end{description}

\subsection{Device-owned --- mission, scheduling, loops, discuss (under \texttt{\small device.*})}

\begin{description}[leftmargin=1.4em,style=nextline,itemsep=0.15em]
  \item[\codew{mission.run}, \codew{mission.think},
  \codew{mission.discuss_round}]
  EAL mission execution exposed as Invoke. Without these an LLM
  inside an agent has no first-class way to compose multi-step /
  multi-agent workflows --- it would have to shell out to the
  \code{easynet mission run} CLI, which forfeits in-process state.
  \codew{mission.discuss_round} drives one round of multi-agent
  conversation; the underlying state shares the
  \code{discuss.*} room model.

  \item[\codew{schedule.add}, \codew{schedule.list},
  \codew{schedule.enable}, \codew{schedule.remove}]
  Cron-style schedule registry over Invoke. The daemon's tick
  runner reads from the same store. Without scheduling-as-an-ability
  every recurring task would need an external cron and would skip
  the AXIOM receipt trail.

  \item[\codew{loop.create}, \codew{loop.status},
  \codew{loop.subscribe}, \codew{loop.cancel}]
  Worker+verify loop primitive. The \code{Live} stream variant
  surfaces per-iteration frames; this is what powers
  ``run-until-passes-verifier'' compositions without polling.

  \item[\codew{discuss.create}, \codew{discuss.post},
  \codew{discuss.list_turns}, \codew{discuss.subscribe}]
  Multi-agent discussion room. The append-only turn log + the
  \code{SnapshotThenLive} stream pattern are how a UI joining
  mid-conversation gets past turns + new ones in one call.
  \codew{mission.discuss_round} writes into this same store.
\end{description}

\subsection{Device-owned --- edge adapters and policy (under \texttt{\small device.*})}

\begin{description}[leftmargin=1.4em,style=nextline,itemsep=0.15em]
  \item[\codew{a2a.bridge.list_skills}, \codew{a2a.bridge.send_task}]
  Inbound A2A: surface this device's local A2A agent cards to
  external A2A clients. Without the bridge, the realm's A2A
  catalogue is invisible to A2A-native clients.

  \item[\codew{a2a.client.send_task}] Outbound A2A: dispatch an
  RPC against a remote A2A agent. The other half of the
  agent-to-agent protocol bridge.

  \item[\codew{mcp.bridge.list_tools}, \codew{mcp.bridge.call_tool}]
  Inbound MCP: project the local ability registry as MCP tools
  for an MCP-native client at the edge. Per ``MCP at the edge,
  not node-to-node'' (RFC §A3) this lives on the device, not
  inside the federation.

  \item[\codew{mcp.client.call}, \codew{mcp.client.list}]
  Outbound MCP: dial a remote MCP server, list and call its
  tools. Pairs with the bridge.

  \item[\codew{ability.access}, admission trust gates]
  Standalone \codew{policy.evaluate}/\codew{policy.simulate} was
  removed from the CLI P0 architecture. Admission is enforced by
  the unified envelope, delegation, trust-anchor, quota, and
  ability-access gates; future permission work must extend that
  path rather than reintroducing a parallel policy ability.
\end{description}

\subsection{Device-owned --- ability publishing primitives (under \texttt{\small device.ability.*})}

\begin{description}[leftmargin=1.4em,style=nextline,itemsep=0.15em]
  \item[\codew{ability.publish}, \codew{ability.unpublish}]
  Publish-layer verbs that push a descriptor through the
  federation advertise pipeline. Without them the only way to
  add a new ability is to redeploy the daemon, which forfeits
  the live-update story.
\end{description}

\subsection{Device-owned --- daemon device-bundle (row 5)}

These are RFC-002 §3.3 abilities the daemon hosts directly, with
no sub-agent wrapper. The catalogue carries exactly one row each.

\begin{description}[leftmargin=1.4em,style=nextline,itemsep=0.15em]
  \item[\codew{device.keyring.create}, \codew{device.keyring.list},
        \codew{device.keyring.get_public},
        \codew{device.keyring.rotate}, \codew{device.keyring.revoke},
        \codew{device.keyring.expire_set},
        \codew{device.keyring.bind_subject},
        \codew{device.keyring.peer_add},
        \codew{device.keyring.peer_list},
        \codew{device.keyring.federate_user_identity_token}]
  RFC-002 keyring: the device-scoped private-key store + signing
  primitive + cross-realm federated identity token issuance. Without
  it, no AXIOM Invoke can be signed and federation breaks at the
  signing-authority check.

  \item[\codew{session.open}] The long-lived bidi RPC session
  the device dials to its realm hub on boot. All
  hub$\to$device forward-Invokes ride this session. Without it the
  hub has no way to push work down to the device.

  \item[\codew{identity.register_pubkey}]
  Pairing-flow primitive: writes the device's Ed25519 public key
  into the hub's realm trust anchor. Without it, the hub's
  admission gate would reject every signed Invoke from the
  device.
\end{description}

% ---------- Sub-agent rows: per-agent (row 4 + 6 + 7) -------------------
\subsection{Agent-owned (per-sub-agent rows: rows 4, 6, 7)}

For every sub-agent registered on the device the catalogue carries
its discover / invoke / chat / api-key triple plus any custom verbs.
On this fixture device the sub-agents are \code{codex} and
\code{web-builder}; with both mounted the rows look like:

\begin{description}[leftmargin=1.4em,style=nextline,itemsep=0.15em]
  \item[\codew{<agent>.chat}] The agent's main entry. Owner =
  \code{agent/<u>.<agent>}. This is what an LLM-routed user
  request hits. Without per-agent chat the only entry would be
  the generic OpenAI adapter (row~9), which is hub-routed --- you
  could not address a specific sub-agent.

  \item[\codew{<agent>.discover}] Per-agent ability catalogue
  filtered to this agent's visibility scope. The LLM uses it to
  enumerate "what verbs may I call as this agent" without
  ever having to read the device-wide catalogue.

  \item[\codew{<agent>.invoke}] Generic Invoke entry-point for
  the agent: ``call any ability you discovered as me''. Without
  it, every per-verb invocation would need its own handler and
  the discover-then-invoke pattern would not compose.

  \item[\codew{<agent>.api_key.create}, \codew{<agent>.api_key.list},
        \codew{<agent>.api_key.revoke}]
  Agent-scoped API key management for the OpenAI-compat surface.
  Each agent owns its own key set. Today's catalogue still shows
  these under the daemon's pages-user identity (\code{self.}); the
  terminal form is per-agent.

  \item[\codew{web-builder.todo_add_task},
        \codew{web-builder.todo_delete_task},
        \codew{web-builder.todo_list_tasks}]
  Custom verbs declared in the \code{web-builder} agent's
  workspace TOMLs. These are the canonical example of row~7
  --- per-agent custom abilities exposed through the same
  registry as the built-ins.
\end{description}

% ---------- Sub-agent meta / skill -------------------
\subsection{Agent-owned (LLM sub-agent surfaces): meta, skill}

These verbs live under \code{web-builder}'s sub-agent in this
fixture --- they are re-exported per agent, not device-wide.

\begin{description}[leftmargin=1.4em,style=nextline,itemsep=0.15em]
  \item[\codew{meta.describe}, \codew{meta.list_abilities},
        \codew{meta.list_resources}]
  Self-introspection over the seven-tuple. \code{describe} is the
  cheap identity+summary surface (federation peer hits this for
  cache checks); \code{list\_abilities} returns the full
  descriptor catalogue; \code{list\_resources} enumerates physical
  /~logical resources (mics, cameras, files) the agent can name as
  subjects. Without these the LLM cannot know its own capabilities.

  \item[\codew{skill.list}, \codew{skill.publish}, \codew{skill.unpublish}]
  Skill (curator-authored ability bundle) lifecycle. Publishing a
  skill mounts a new ability under the agent's namespace at runtime;
  this is how the \code{alive-video}, \code{macos}, etc. skills land
  on a device without redeploying.

\end{description}

% ---------- Hub voice authority (10 rows) -------------------
\subsection{Hub-owned --- voice signaling and model media (row 9)}

Microphone and speaker handles remain Device-owned resources under
\codew{mic.*} and \codew{speaker.*}.  The \code{voice.*} namespace is
realm-shared state and therefore has one canonical owner:
\codew{easynet:///r/<r>/hub}.  The call registry is keyed by the pair
\code{(hub authority URA, call\_id)}; Device and LLM profiles must not
re-export these descriptors, and there is no Device proxy alias.

\begin{description}[leftmargin=1.4em,style=nextline,itemsep=0.15em]
  \item[\codew{voice.create_call}, \codew{voice.join_call},
        \codew{voice.leave_call}, \codew{voice.end_call},
        \codew{voice.list_calls}, \codew{voice.show_call},
        \codew{voice.watch_call}, \codew{voice.report_metrics}]
  Realm call signaling and participant state.  Every handler validates
  the Invocation callee as the Hub authority before reading or mutating
  the shared call registry.

  \item[\codew{voice.subscribe}, \codew{voice.transcribe}]
  Hub-governed TTS and ASR media planes.  Their subjects name voice or
  ASR model resources; physical capture and playback remain separate
  Device-owned microphone and speaker Invocations.
\end{description}

% ---------- consent helper (3 rows) -----------------------------------
\subsection{Agent-owned --- consent default helper (row 8 sub-case~a)}

Owner URA: \codew{easynet:///r/easynet.run/agent/test.consent-default-0}.

\begin{description}[leftmargin=1.4em,style=nextline,itemsep=0.15em]
  \item[\codew{consent.subscribe}] Subscribe to the
  approval-broker pending queue. \code{SnapshotThenLive} stream;
  a UI joining mid-flight gets currently-pending requests + new
  ones.
  \item[\codew{consent.list_pending}] One-shot snapshot of the
  same queue (the unary sibling).
  \item[\codew{consent.decide}] Resolve a pending permission
  request by id; decision $\in$~\{\code{allow}, \code{deny},
  \code{deny\_with\_reason}\}.
\end{description}

\paragraph{Why these need a hosted helper agent (not just a daemon adapter).}
Consent decisions carry user-attributable signatures: the \code{deny}
on a destructive Invoke is itself an Invoke, signed by the user (via
the consent helper agent). A daemon-side adapter cannot sign as the
user --- only an agent the user owns can. Hence the helper-agent
ownership form (row~8 sub-case~a).

% =====================================================================
\section*{References}

\begin{itemize}[leftmargin=1.6em,itemsep=0.1em]
  \item \codew{docs/rfc/AXON-RFC-001-plan-v4.1.5.md} --- URA grammar §A.URA, AXIOM seven-tuple §9, ability registry §18.
  \item \codew{docs/rfc/AXON-RFC-002-session-provider.md} --- §3.3 device-bundle keyring (row~5).
  \item \codew{docs/rfc/AXON-RFC-006-stateful-easynet.pdf} --- §C OpenAI compatibility surface (row~9).
  \item \codew{src/core/ura/mod.rs} --- canonical URA builders + parser; the §1 grammar is reflected here verbatim.
  \item \codew{src/daemon/ability/descriptors/surface.rs} --- \code{AbilityDescriptor.owner\_ura}; v4.1.5 §9 \code{callee} constraint comment.
  \item \codew{src/daemon/ability/builtins/governance/meta.rs} --- \code{list\_abilities\_handler} synth path.
  \item \codew{src/cli/abilities.rs} --- \code{easynet ability list} render layer.
\end{itemize}

\vspace{1em}
\hrule
\vspace{0.4em}
\begin{flushright}\small
  \emph{End of specification.}\\
  Updates require an RFC bump.
\end{flushright}

\end{document}
